\section{Context Mining}
\label{mining_method}

The foundation of \sysname{} is its context mining engine. This component is tasked with systematically identifying and extracting information that may be relevant to analyzing a specific line of configuration. The core challenge lies in the complex nature of `relevance' itself, which is not defined by mere proximity in a file. As discussed in Section~\ref{sec:motivation}, configurations contain numerous connections within and across devices. 
% Therefore, to emulate an expert's intuition and capture the full picture of a configuration's correctness—including its local interactions, \aaron{What are ``local interactions?''} network-wide patterns, and distant dependencies—we design a mining process that is guided by configurations' hierarchical structure, relative uniformity~\cite{le2006minerals, kakarla2020finding}, and abundance of references~\cite{benson2009complexitymetrics}.
Therefore, to emulate an expert's intuition, we design a mining process guided by configurations' hierarchical structure, relative uniformity~\cite{le2006minerals, kakarla2020finding}, and abundance of references~\cite{benson2009complexitymetrics} in order to capture the full picture of a configuration's correctness by revealing its local interactions (\ie, how settings within a single block define its behavior), its adherence to network-wide patterns (the implicit design templates), and its distant dependencies (the explicit links between a statement and its definition elsewhere).


\subsection{Representing configurations}
\label{sec:mining_method:structure}

%\begin{wrapfigure}{r}{0.5\columnwidth}
%    \centering
%    \fbox{\includegraphics[width=0.48\columnwidth]{figs/tree_form.png}}
%    \caption{Example snippet of tree-formatted Juniper router configuration file. \aaron{Groups are a Juniper-specific feature; can we show a snippet of the configuration that is more common across vendors (e.g., interfaces, routing protocols, filters)?}}
%    \label{fig:tree}
%\end{wrapfigure}

Since configurations have a structured, hierarchical format (discussed in Section~\ref{sec:motivation}), we represent each device's configuration as a tree ($T$).
%Configurations are typically written in a structured, hierarchical format that can be effectively modeled as a tree (\(T\)), as shown in the example in Figure~\ref{fig:tree}. 
Each node (\(V\)) corresponds to a specific 
%configuration element, with child nodes providing further specificity.
category (\eg, `{\tt interfaces}', `{\tt filter}'), instance (\eg, `{\tt ge-1/0/1}', `{\tt EXAMPLE-FILTER}'), or parameter (\eg, `{\tt mtu}', `{\tt input}') within the configuration.
%and edges convey the nesting of 
%represent the relationships/dependencies between
%these elements.
A complete path \(
P = \{ \top \rightarrow V_1 \rightarrow V_2 \rightarrow \dots \rightarrow V_k = v_k \}
\) from the root ($\top$) to a leaf ($V_k$) represents a single, \textit{fully-qualified configuration statement} (FQCS). For example, the path listed at the top of Figure~\ref{fig:context_mine} corresponds to line 7 in Figure~\ref{fig:example_config}, where the leaf node ($V_k = v_k$) represents the configuration parameter (`{\tt input}') and its associated value (`{\tt EXAMPLE-FILTER}'), and the preceding nodes ($V_1, V_2, \dots, V_{k-1}$) represent the configuration sections the parameter is nested within (`{\tt interfaces}', `{\tt ge-0/0/1}', ..., `{\tt filter}').

%Let any unique path \(
%P = \{ \top \rightarrow V_1 \rightarrow V_2 \rightarrow \dots \rightarrow V_k = v_k \}
%\) in this tree have \(k\) nodes. The root node (\(\top\)) represents the entry point of the configuration file. The immediate following node (\(V_1\)) corresponds to the broadest configuration category in this path. Intermediate nodes (\( V_2, \dots, V_{k-1} \)) represent increasingly nested configuration sections. Finally, the leaf node \(V_k\) represents a configurable parameter with its associated value \(v_k\). A complete path from the root to a leaf therefore represents a single, \textit{fully-qualified configuration statement} (FQCS), which may be composed of multiple hierarchical lines in the raw text (\eg, the setting `allowed vlans 123` is found within the `interface GE1/1` stanza). %


%\subsection{Contextual Relevance from Different Aspects:} 
\subsection{Principal configuration contexts} 
Based on the characteristics of configurations discussed in Section~\ref{sec:motivation}, we identify three principal forms of context to extract from configurations: (1) neighboring configuration, (2) similar configuration, and (3) referenced configuration.
We define each type of context relative to a specific FQCS ($P = \{ \top \rightarrow V_1 \rightarrow \dots \rightarrow V_k = v_k \}$) under review. As demonstrated by the examples in Figure~\ref{fig:context_mine}, each type of context provides unique insights that contribute to the evaluation of a FQCS. 


% \begin{figure*}[t]
%     \centering
%     \begin{mdframed}
%     \begin{minipage}{0.5\textwidth}
%     \footnotesize % Dynamically scales down font
%     \mypara{Example: Context Mining for a Specific Configuration Line}
    
%     \textit{Configuration Line Under Review:}
%     \begin{Verbatim}
% groups -> INTERFACE-BACKBONE -> interfaces -> <ge-*> -> unit <*>
% -> family inet -> filter -> input-list = "inet-sample"
%     \end{Verbatim}
%     \hrulefill

%     \mypara{Extracted Neighboring Configurations}
%     \begin{Verbatim}
% (a) groups -> INTERFACE-BACKBONE -> interfaces -> <ge-*> -> unit <*>
%     -> family inet -> filter -> input-list = "backbone-in"
% (b) groups -> INTERFACE-BACKBONE -> interfaces -> <ge-*> -> unit <*>
%     -> family inet -> filter -> output = "interface-out"
%     \end{Verbatim}

%     \mypara{Extracted Similar Configurations}
%     \begin{Verbatim}
% (a) groups -> INTERFACE-BACKBONE -> interfaces -> <ge-*> -> unit <*>
%     -> family mpls -> filter -> input-list = "mpls-sample"
%     \end{Verbatim}

%     \mypara{Extracted Referenced Configurations}

%     \mysubpara{(1) Intra-Device Reference}
%     \begin{Verbatim}
% (a) groups -> firewall -> family inet -> filter inet-sample -> term default
%     -> then = sample
% (b) groups -> firewall -> family inet -> filter inet-sample -> term default
%     -> then = next term
%     \end{Verbatim}

%     \mysubpara{(2) Cross-Device Reference}
%     \textit{New Example Line Under Review (on Router A, IP 10.1.2.2):}
%     \begin{Verbatim}
% protocols -> bgp -> group "PEERS" -> neighbor 10.1.2.3 -> description = "to-router-B"
%     \end{Verbatim}
%     \textit{Extracted Context (from Router B's configuration file):}
%     \begin{Verbatim}
% (a) protocols -> bgp -> group "PEERS" -> neighbor 10.1.2.2 -> description = "to-router-A"
%     \end{Verbatim}
%     \end{minipage}
%     \end{mdframed}

%     \caption{An illustration of the four types of context mined by \sysname{} for a given configuration line. The system extracts neighboring, similar, and both intra- and cross-device referenced configurations. \aaron{This needs to be updated to match the example configuration in Figure~\ref{fig:example_config}.}}
%     \label{fig:context_mine}
% \end{figure*}


\begin{figure*}[t]
    \centering
    \begin{mdframed}
    \begin{minipage}{0.95\textwidth} % Using a slightly wider minipage for clarity
    \footnotesize % Dynamically scales down font
    
    \mypara{Example: Context Mining Based on Figure~\ref{fig:example_config}}
    
    \vspace{1em}
    \textit{Configuration Line Under Review (from line 7):}
    \begin{Verbatim}
interfaces -> ge-0/0/1 -> unit 0 -> family inet -> filter -> input = "EXAMPLE-FILTER"
    \end{Verbatim}
    \hrulefill

    \mypara{Extracted Neighboring Configurations}
    \textit{(Other settings for the same interface unit)}
    \begin{Verbatim}
(a) interfaces -> ge-0/0/1 -> unit 0 -> family inet -> address = "10.0.1.1/24"
    \end{Verbatim}

    \mypara{Extracted Similar Configurations}
    \textit{(The same parameter used on a different, but similar, interface)}
    \begin{Verbatim}
(a) interfaces -> ge-0/0/2 -> unit 0 -> family inet -> filter -> input = "EXAMPLE-FILTER"
    \end{Verbatim}

    \mypara{Extracted Referenced Configurations}

    \mysubpara{(1) Intra-Device Reference}
    \textit{(Finding the definition of the referenced "EXAMPLE-FILTER")}
    \begin{Verbatim}
(a) firewall -> family inet -> filter EXAMPLE-FILTER -> term ALLOW-INTERNAL 
    -> from -> source-address = "10.0.0.0/16"
(b) firewall -> family inet -> filter EXAMPLE-FILTER -> term ALLOW-INTERNAL 
    -> then -> accept
    \end{Verbatim}

    \mysubpara{(2) Cross-Device Reference}
    \textit{(New Example Line Under Review (from line 24) on Router A (10.0.2.1)):}
    \begin{Verbatim}
protocols -> bgp -> group "EBGP-PEER" -> neighbor = "10.0.2.2"
    \end{Verbatim}
    \textit{Extracted Context (from peer Router B's configuration file):}
    \begin{Verbatim}
(a) protocols -> bgp -> group "EBGP-PEER" -> neighbor = "10.0.2.1"
    \end{Verbatim}
    
    \end{minipage}
    \end{mdframed}

    \caption{An illustration of the types of context mined by \sysname{}, with concrete examples drawn from the sample configuration in Figure~\ref{fig:example_config}. The system extracts neighboring, similar, and both intra- and cross-device referenced configurations.}
    \label{fig:context_mine}
\end{figure*}

\mypara{Neighboring configurations ($N_m(P)$):}
% Neighboring configurations are lines that are structurally close to the FQCS under review, typically within the same functional block. They provide immediate, local context, revealing potential dependencies or conflicts. \aaron{It's a little confusing to talk about dependencies and conflicts here, because dependencies seems more connected to referenced configurations and conflicts seems more connected to similar configurations. Is there another way we can motivate neighboring configurations?} For instance, when analyzing an IP address on an interface, other settings for that same interface (like its MTU or ACLs) are crucial neighboring context.
Neighboring configurations are lines that are structurally close to the FQCS under review, typically within the same functional block. They provide the immediate, local context necessary to understand the holistic state and intended behavior of a specific entity, such as an interface or routing protocol. For instance, when analyzing an IP address on an interface, its neighboring context—such as whether the interface is administratively shut down, its assigned MTU, or the access lists applied to it—is crucial for correctly interpreting the operational impact of that IP address.


Formally, we define neighboring configurations as the set of paths $P'$ that share a common ancestry with $P$ up to a certain depth, $m$. Let $V_i(P)$ be the $i$-th node in path $P$. The set of neighboring paths is:
\[
N_m(P) = \{ P' \in T \mid V_i(P') = V_i(P) \text{ for } 1 \le i \le m, \text{ and } P' \neq P \}
\]
%This corrected definition ensures the shared nodes are identical. 

The depth parameter $m$ controls the scope of the context. A larger $m$ results in a tighter, more relevant neighborhood.
%\textit{Design Justification:} 
A simple and effective choice is to set $m = k-1$, meaning we extract all sibling configurations that share the same direct parent. This captures the most immediate context without including potentially irrelevant data from higher up in the hierarchy. While we use this default for our evaluation, we acknowledge that a full study on the impact of different values of $m$ is a valuable area for future work.

\mypara{Similar configurations ($S(P)$):}
Similar configurations involve the same type of parameter but applied in slightly different contexts (\eg, the MTU settings on all other interfaces across the network). Comparing these helps detect inconsistencies or deviations from standard practices---which, as discussed in Section~\ref{sec:motivation:llms}, may or may not be intentional.

% Formally, we define similar configurations as the set of paths $P'$ that share the same broadest category node ($V_1$) and the same final parameter node ($V_k$), but differ in their intermediate nodes.
% \[
% S(P) = \{ P' \in T \mid V_1(P')=V_1(P), V_k(P')=V_k(P), P' \neq P \}
% \]
% For example, this allows us to find all `MTU' settings ($V_k$) under the `interfaces' category ($V_1$) across the entire device.

% We require a match up to the $V_1$ node because it typically represents a major functional block (like `interfaces' or `routing-policies'). This provides a meaningful balance; matching only on the root node ($\top$) would be too broad, while requiring deeper matches (like $V_2$) would be too restrictive and might miss relevant comparisons across different sub-sections. \aaron{Should we include an example of why these alternatives would be too broad or restrictive: e.g., finding similar uses of the `action' parameter under $\top$ would include both route filters and packet filters, while matching by a specific interface would find no similar configurations?} \aaron{Should we include a depth parameter in $S(P)$ similar to $N_m(P)$? As a motivating example: `routing' is a $V_1$ node in Juniper router configurations, with specific protocols (`bgp`, `ospf`) as $V_2$ nodes;  it seems odd to compare BGP settings with OSPF settings, which occur when only matching on $V_1$. At the same time, it is unlikely that there will be overlap in $V_K$ between BGP and OSPF settings, so maybe this flexibility is irrelevant. Should we be conducting some empirical analysis of configuration languages to motivate the choice of matching only on $V_1$?}


To provide flexibility, we define similar configurations using a shared ancestry depth, $m$, similar to our approach for neighboring configurations. Formally, we define the set of similar paths as those that share the same ancestry up to depth $m$ and terminate with the same parameter node ($V_k$), but differ in the path between $m$ and $k$.
\[
S_m(P) = \{ P' \in T \mid V_i(P')=V_i(P) \text{ for } 1 \le i \le m, V_k(P')=V_k(P), \text{ and } P' \neq P \}
\]

The choice of the depth parameter $m$ allows for a crucial trade-off between the breadth and relevance of the comparison. Our default choice is to set $m=1$, which typically represents a major functional block (like `interfaces' or `protocols'). This provides a meaningful balance, which can be illustrated by considering the alternatives:
\begin{itemize}
    \item A choice of $m=0$ (matching only on the root node $\top$) would be too broad. For example, searching for a parameter like `action' would incorrectly group together firewall filter actions and route policy actions, which are semantically very different.
    \item A choice of $m$ that is too high (\eg, $m = k-2$, matching on a specific interface like `ge-0/0/1') would be too restrictive. A search for similar MTU settings would find no configurations on other interfaces, defeating the purpose of the comparison.
\end{itemize}

The flexibility of $m$ also helps manage comparisons within large functional blocks. For instance, in a Juniper configuration, `protocols' is a $V_1$ node containing specific protocols like `bgp' and `ospf' as $V_2$ nodes. Our default of $m=1$ would compare settings across BGP and OSPF. This may seem odd, but in practice, the final parameter names ($V_k$) are typically distinct between protocols, so irrelevant comparisons are rare. However, for a more targeted analysis, a user could set $m=2$ to find similar settings only within BGP or within OSPF. While our default provides a reasonable starting point, to suit different analysis cases, \sysname{} supports the flexibility to choose optimal default values for $m$ for different parameter types.



\mypara{Referenced Configurations (Intra- and Cross-Device):}
Referenced configurations provide the definitions for user-defined identifiers --- custom labels created by an administrator to name a specific object or policy, such as an access-list name, a policy-map, or a BGP peer group. Referenced configurations are the stanzas that provide the actual definitions for these identifiers. Mining this context is crucial for understanding dependencies, but the nature of these references can be either internal to a single device's configuration or external, pointing to another device on the network. \sysname{} is designed to handle both.

\mysubpara{1. Intra-Device Referencing:}
This is the most common form, where a value in one part of a configuration refers to a policy or object defined elsewhere in the same file. For example, a line like `import-policy = PolicyX` is a reference. To resolve this, the context mining engine searches for the stanza where `PolicyX` is defined, such as the `policy-options $\rightarrow$ policy-statement PolicyX $\rightarrow$ ...' block, which contains the actual rules for that policy.

Formally, for a given path $P$ in a device's configuration tree $T$, we define this set of intra-device references, $R_{intra}(P)$, as the set of paths $P'$ where the value of the original path, $v_k(P)$, appears as an intermediate node.
\[
R_{intra}(P) = \{ P' \in T \mid P' =  \top \rightarrow \dots \rightarrow v_k(P) \rightarrow \dots \rightarrow V_j = v_j \}
\]
Here, $v_k(P)$ is no longer a leaf node but part of the path in $P'$, indicating that it is being defined or elaborated upon.


% \mypara{Differentiating Pre-defined Keywords from User-defined Identifiers:} \aaron{Still need to refine and integrate this part}

A primary challenge in extracting intra-device referenceable configurations is the risk of incorporating irrelevant context. This risk stems from the difficulty of distinguishing between \textit{pre-defined parameter values} and \textit{custom user-defined parameter values}.
Pre-defined keywords are part of the vendor's standard configuration language and have a fixed, universal meaning (\eg, the boolean values `true' and `false', interface modes like `trunk' or `access', or protocol names like `any' or `ospf'). In contrast, user-defined identifiers are custom labels created by an administrator to name a specific object or policy (\eg, a policy map name like `QOS-VOIP-POLICY', an access list name such as `BLOCK\_GUEST\_ACCESS', or a VLAN ID like `1000'). This distinction is critical because it dictates our context-mining strategy: a user-defined identifier acts as a pointer, implying that its full definition exists elsewhere in the configuration and \textit{must} be retrieved. 
Misidentifying a custom policy name as a pre-defined keyword would prevent the system from searching for its essential definition, likely causing it to miss critical dependency-related errors. Conversely, mistakenly treating a standard keyword like `true' as user-defined would trigger a pointless and computationally expensive search for a non-existent or mismatched definition.



To solve this, we initially considered a simple heuristic based on existence checks in the configuration's hierarchical tree structure. This method assumed that any parameter value that also appeared as a parent or intermediate node elsewhere in the configuration must be user-defined. However, this approach proved brittle, as it could be fooled by coincidences in naming and lacked a true understanding of the configuration's grammar. We then explored a more advanced majority-voting mechanism, which attempted to enforce consistency by assuming that a given parameter type (like `timeout') should always take the same class of value (pre-defined or user-defined). This, too, had a critical flaw: it could not distinguish between identical parameter values used in different contexts. For example, it would struggle to differentiate between the user-defined `vlan-id 1000' and the user-defined `access-list 1000', as it would incorrectly try to enforce a uniform classification on the value `1000' based on its other uses, imposing a false consistency.


Therefore, \sysname{}'s final design discards these rule-based heuristics in favor of a more direct and robust \textit{LLM-based approach}. As part of its initial query, our framework explicitly instructs the LLM to classify the value not in isolation, but \textit{in the context of the full configuration line it appears in}. By leveraging its vast training on vendor-specific documentation and millions of real-world configuration examples, the LLM has learned the expected syntax and value patterns for different configuration parameters. It understands that the value following a parameter like `vlan-id' is a user-created label, while the value following a command like `ip routing' is a pre-defined boolean keyword. This allows it to easily and accurately determine that `1000' is a distinct, user-defined identifier in the context of a VLAN versus an ACL, correctly handling the ambiguity that plagued the previous methods. This approach requires no complex, hand-crafted rules and has proven to be significantly more accurate and reliable.

\mysubpara{2. Cross-Device Referencing:}
A more complex challenge arises when a user-defined identifier refers to another device on the network. This is common in configurations for routing protocols, VPNs, or services. For example, in a BGP configuration, the line `neighbor 10.1.1.2 remote-as 65001' contains an identifier, `10.1.1.2', for a peer router. The most critical context for verifying this line is not on the current device, but on the device at `10.1.1.2'.

Our methodology for this is a three-step process:
\begin{enumerate}
    \item \textbf{Identify External Reference:} The system first identifies that a value (like an IP address or hostname) is an identifier for an external device.
    \item \textbf{Locate Target Device:} It then looks up the corresponding configuration file for that target device from our set of network-wide configurations.
    \item \textbf{Find Reciprocal Configuration:} Finally, it searches within the target device's configuration for lines that reference the \textit{original} device back. For the BGP example, if our original device's IP is `10.1.1.1`, the system would search the configuration of `10.1.1.2` for the corresponding `neighbor 10.1.1.1 ...` statement.
\end{enumerate}

Formally, let the current device's configuration tree be $T_A$ and its network identifier be $id(T_A)$. If the value $v_k(P)$ in a path $P \in T_A$ is an identifier for another device with configuration tree $T_B$, the set of cross-device references, $R_{cross}(P)$, is defined as:
\[
R_{cross}(P) = \{ P' \in T_B \mid id(T_A) \in P' \}
\]
This captures all configuration statements on the target device that mention the original device. This process is essential for verifying connectivity and detecting symmetric misconfigurations where, for example, a session is configured on one side but not the other.
The complete set of referenced configurations, $R(P)$, encompasses both intra-device definitions found within the same configuration file and cross-device configurations located on a target device that refer back to the original device.

\mypara{Composable Context Primitives:}
It is important to note that the three principal context types---Neighboring, Similar, and Referenced---serve as atomic building blocks, or \textit{primitives}. A key strength of our framework is the ability to compose these primitives to form more complex and powerful context queries. For example, we can compose S(R(P)), which translates to finding configurations that are \textit{similar} to the \textit{definition} of a value in the line under review. Conversely, a composition of R(S(P)) would find the \textit{definitions} of values that appear in \textit{similar} configurations. This compositional nature, which also allows for queries like R(N(P)) (the definitions of values in neighboring lines), gives \sysname{} the flexibility to emulate the multi-step, intuitive reasoning of a human expert.